\documentclass[11pt,a4paper]{article}
\usepackage[margin=2.3cm]{geometry}
\usepackage[T1]{fontenc}
\usepackage{amsmath,amssymb}
\usepackage{graphicx}
\usepackage{booktabs}
\usepackage{xcolor}
\usepackage{hyperref}
\usepackage{listings}
\usepackage{fancyhdr}
\usepackage{float}
\definecolor{upesnavy}{HTML}{1F3864}
\pagestyle{fancy}
\fancyhf{}
\lhead{\small Asymptotic Notations}
\rhead{\small Yashvardhan Singh}
\cfoot{\thepage}
\lstset{basicstyle=\ttfamily\small,breaklines=true,frame=single,numbers=left}
\newcommand{\studentname}{Yashvardhan Singh}
\newcommand{\sapid}{590021777}
\newcommand{\batch}{B17}
\newcommand{\topic}{Asymptotic Notations: Big O, Big Omega, Big Theta}
\newcommand{\submissiondate}{15 October 2026}
\begin{document}
\begin{titlepage}
\centering
{\Large UPES, Dehradun\par}
{\large School of Computer Science \par}
{\large Artificial Intelligence Cluster\par}
\vspace{2cm}
{\color{upesnavy}\rule{\linewidth}{1.2pt}\par}
\vspace{0.5cm}
{\Huge\bfseries \topic\par}
\vspace{0.5cm}
{\Large Mid-Semester Assignment Report\par}
\vspace{0.5cm}
{\color{upesnavy}\rule{\linewidth}{1.2pt}\par}
\vspace{2cm}
\begin{tabular}{rl}
\textbf{Course:} & Design and Analysis of Algorithms (CSEG3053)\\[0.3cm]
\textbf{Student:} & \studentname\\[0.3cm]
\textbf{SAP ID:} & \sapid\\[0.3cm]
\textbf{Batch:} & \batch\\[0.3cm]
\textbf{Semester:} & III (Aug--Dec 2026)\\[0.3cm]
\textbf{Submitted to:} & Kunwar Siddharth, Assistant Professor\\[0.3cm]
\textbf{Date:} & \submissiondate
\end{tabular}
\vfill
\end{titlepage}

\begin{abstract}
Asymptotic notation is a mathematical language used to describe how an algorithm's resource requirements grow as the input becomes large. This report explains three central notations: Big O, Big Omega, and Big Theta. Big O gives an asymptotic upper bound, Big Omega gives a lower bound, and Big Theta gives a tight bound when upper and lower bounds have the same order of growth. The report introduces the idea using simple examples, states formal definitions, and demonstrates the analysis of a loop-based algorithm. A Python animation plots constant, logarithmic, linear, linearithmic, and quadratic growth as the input size increases. The key lesson is that asymptotic analysis helps compare algorithm scalability independently of a particular computer, while constants and lower-order terms still matter for small inputs and practical performance.
\end{abstract}

\section{Introduction}
An algorithm is a sequence of steps used to solve a problem. When choosing between two algorithms, it is useful to know how their running time or memory use changes as the input grows. Measuring seconds on one laptop is not enough: the result depends on hardware, programming language, implementation, and current system load. Asymptotic analysis instead studies the growth trend as input size $n$ becomes large.

This topic is part of Design and Analysis of Algorithms. It is useful when selecting searching, sorting, graph, and data-processing methods. For example, an algorithm that performs roughly $n^2$ operations may be acceptable for 20 items but much slower than an $n\log n$ algorithm for 100,000 items. Asymptotic notation gives a concise way to describe and compare such trends.

\section{Intuition and Core Idea}
Big O describes an upper bound, Big Omega describes a lower bound, and Big Theta describes a tight bound. They describe mathematical bounds on a function; they are not automatically synonyms for worst, best, and average cases. Those cases describe which input or scenario is being measured, while asymptotic notation describes the growth bound of the resulting function.

\begin{center}
\fcolorbox{upesnavy}{white}{\parbox{0.88\linewidth}{
\textbf{In simple words:}\\
$O(g(n))$: the function eventually grows no faster than a constant multiple of $g(n)$.\\
$\Omega(g(n))$: the function eventually grows at least as fast as a positive constant multiple of $g(n)$.\\
$\Theta(g(n))$: the function is bounded above and below by constant multiples of $g(n)$.
}}
\end{center}

\begin{figure}[H]
\centering
\includegraphics[width=0.92\linewidth]{figures/growth_rates.png}
\caption{Illustrative growth rates. The plotted values compare trends, not actual timings from a particular computer.}
\end{figure}

\section{Mathematical Formulation}
Let $f(n)$ be a non-negative function representing an algorithm's operation count for sufficiently large input size $n$.

\subsection{Big O: asymptotic upper bound}
$f(n)=O(g(n))$ if there exist positive constants $c$ and $n_0$ such that
\[
0\le f(n)\le c\,g(n)\quad\text{for all }n\ge n_0.
\]
Example: $f(n)=3n^2+5n+2$ is $O(n^2)$ because for sufficiently large $n$, the lower-order terms can be bounded by a constant multiple of $n^2$.

\subsection{Big Omega: asymptotic lower bound}
$f(n)=\Omega(g(n))$ if there exist positive constants $c$ and $n_0$ such that
\[
0\le c\,g(n)\le f(n)\quad\text{for all }n\ge n_0.
\]
The same polynomial $3n^2+5n+2$ is $\Omega(n^2)$ because it is always at least $3n^2$ for positive $n$.

\subsection{Big Theta: tight asymptotic bound}
$f(n)=\Theta(g(n))$ if there exist positive constants $c_1,c_2,n_0$ such that
\[
0\le c_1g(n)\le f(n)\le c_2g(n)\quad\text{for all }n\ge n_0.
\]
Since $3n^2+5n+2$ is both $O(n^2)$ and $\Omega(n^2)$, it is $\Theta(n^2)$.

\section{Algorithm / Procedure}
The following pseudocode counts the number of times an inner statement runs. The outer loop executes $n$ times and the inner loop executes $n$ times for each outer iteration, so the total count is $n\times n=n^2$.

\begin{lstlisting}[language=Python,caption={Counting operations in a nested loop}]
def count_operations(n):
    count = 0
    for i in range(n):
        for j in range(n):
            count += 1
    return count
\end{lstlisting}

For a positive input $n$, the function returns $n^2$. Therefore its operation count is $\Theta(n^2)$, and consequently it is also $O(n^2)$ and $\Omega(n^2)$. The notation $\Theta(n^2)$ is the most informative of these three descriptions because it gives a tight bound.

\section{Worked Example (Step by Step)}
Consider $f(n)=3n^2+5n+2$.

\begin{enumerate}
\item Identify the terms: $3n^2$, $5n$, and $2$.
\item For large $n$, the quadratic term grows faster than the linear and constant terms.
\item For $n\ge 1$, $n\le n^2$ and $1\le n^2$, so
\[
3n^2+5n+2\le 3n^2+5n^2+2n^2=10n^2.
\]
Thus $f(n)\le 10n^2$, proving $f(n)=O(n^2)$ with $c=10,n_0=1$.
\item Also $f(n)\ge 3n^2$ for $n\ge1$, proving $f(n)=\Omega(n^2)$ with $c=3,n_0=1$.
\item Combining both bounds gives $f(n)=\Theta(n^2)$.
\end{enumerate}

\begin{center}
\begin{tabular}{rrr}
\toprule
$n$ & $f(n)=3n^2+5n+2$ & $n^2$\\
\midrule
1 & 10 & 1\\
2 & 24 & 4\\
5 & 102 & 25\\
10 & 352 & 100\\
\bottomrule
\end{tabular}
\end{center}
The values are illustrative operation counts. Big O and related notations describe growth for sufficiently large inputs, not exact equality between functions.

\section{Animation / Software Demonstration}
The demonstration is implemented in Python using NumPy and Matplotlib. It animates the increase in input size and compares five functions: $1$, $\log_2 n$, $n$, $n\log_2 n$, and $n^2$. The animation makes it visible that quadratic growth eventually rises much faster than the other plotted functions.

\textbf{Files:} \texttt{main.py}, \texttt{requirements.txt}, \texttt{growth\_rates\_animation.gif}.\\
\textbf{How to run:} Install Python 3, run \texttt{pip install -r requirements.txt}, then run \texttt{python main.py}. The program displays the animation and saves a GIF in the current directory. The supplied GIF can also be opened directly without running the code.

\begin{figure}[H]\centering
\includegraphics[width=0.75\linewidth]{figures/frame1.png}
\caption{Keyframe 1: small input size.}\end{figure}
\begin{figure}[H]\centering
\includegraphics[width=0.75\linewidth]{figures/frame2.png}
\caption{Keyframe 2: early growth.}\end{figure}
\begin{figure}[H]\centering
\includegraphics[width=0.75\linewidth]{figures/frame3.png}
\caption{Keyframe 3: increasing separation.}\end{figure}
\begin{figure}[H]\centering
\includegraphics[width=0.75\linewidth]{figures/frame4.png}
\caption{Keyframe 4: larger input size.}\end{figure}

\section{Analysis}
\begin{center}
\begin{tabular}{p{2.5cm}p{3.2cm}p{3.2cm}}
\toprule
\textbf{Criterion} & \textbf{Asymptotic notation} & \textbf{Exact benchmarking}\\
\midrule
Purpose & Describes growth trend & Measures actual execution time\\
Strength & Hardware-independent comparison & Reflects a real implementation\\
Limitation & Hides constants and lower-order details & Depends on hardware and environment\\
Best use & Predicting scalability & Testing practical performance\\
\bottomrule
\end{tabular}
\end{center}

Typical growth classes, from slower to faster for sufficiently large $n$, are $O(1)$, $O(\log n)$, $O(n)$, $O(n\log n)$, and $O(n^2)$. This ordering helps explain why efficient sorting algorithms such as merge sort, with $\Theta(n\log n)$ time, usually scale better than simple quadratic sorting methods for large inputs. However, constants, memory use, and input characteristics can affect practical choices.

\section{Applications}
\begin{enumerate}
\item \textbf{Searching:} Linear search has $O(n)$ worst-case time, while binary search on sorted data has $O(\log n)$ time.
\item \textbf{Sorting:} Merge sort has $\Theta(n\log n)$ time, while simple nested-loop sorting approaches commonly have $\Theta(n^2)$ time.
\item \textbf{Data processing:} Complexity analysis helps estimate how processing costs change as datasets become larger.
\end{enumerate}

\section{Conclusion}
This report introduced Big O, Big Omega, and Big Theta as upper, lower, and tight asymptotic bounds. A polynomial example showed how to prove a tight bound, while a nested-loop program demonstrated quadratic operation growth. The animation compared common growth rates visually. The main conclusion is that asymptotic analysis is a useful tool for reasoning about scalability, but it should be combined with practical testing when choosing an implementation.

\section*{Declaration of AI Tool Use}
\textbf{Tools used:} ChatGPT assisted with structuring and drafting explanations, and with generating example code and figures.\\
\textbf{What for:} Initial drafting, explanation of asymptotic bounds, and code/figure preparation.\\
\textbf{Approximate share of AI-generated content:} Substantial; review and revise the report to match your course requirements and your own work before submission.\\
I confirm that I have reviewed the report and will ensure that the submitted work follows the course policy.\\[1cm]
Signature: \rule{4cm}{0.4pt}\hfill Date: \rule{3cm}{0.4pt}

\begin{thebibliography}{9}
\bibitem{cormen}
Thomas H. Cormen, Charles E. Leiserson, Ronald L. Rivest, and Clifford Stein.
\textit{Introduction to Algorithms}, 4th ed. MIT Press, 2022.
\bibitem{knuth}
Donald E. Knuth. \textit{The Art of Computer Programming, Volume 1: Fundamental Algorithms}, 3rd ed. Addison-Wesley, 1997.
\bibitem{skiena}
Steven S. Skiena. \textit{The Algorithm Design Manual}, 3rd ed. Springer, 2020.
\end{thebibliography}
\end{document}
